\section{Sujet n\textsuperscript{o}\,15}
\noindent\begin{minipage}{0.5\linewidth}\footnotesize Office du Baccalauréat du Cameroun\end{minipage}%
\hfill\begin{minipage}{0.45\linewidth}\footnotesize\raggedleft Examen : \textbf{Baccalauréat} · Série : \textbf{C/E}\\
Épreuve : \textbf{Mathématiques} · Durée : \textbf{4\,h} · Coef. : \textbf{7}\end{minipage}
\par\smallskip{\color{onbo}\rule{\linewidth}{1pt}}

\subsection*{Partie A — Évaluation des ressources \hfill (15 points)}

\noindent\textbf{Exercice 1.} \hfill\textit{(3 points)}\\
Une microfinance de Douala étudie le lien entre le nombre $x$ d'années d'ancienneté d'un
client et le montant $y$ (en centaines de milliers de francs CFA) qu'il a épargné. Sur un
échantillon de $5$ clients, on relève le tableau suivant :
\begin{center}\renewcommand{\arraystretch}{1.2}
\begin{tabular}{|c|c|c|c|c|c|}\hline
$x_i$ (années) & $1$ & $2$ & $3$ & $4$ & $5$ \\\hline
$y_i$ (centaines de milliers) & $7$ & $11$ & $14$ & $18$ & $20$ \\\hline
\end{tabular}\end{center}
\begin{enumerate}[label=\arabic*)]
\item Calculer les moyennes $\overline{x}$ et $\overline{y}$, ainsi que la variance
$V(x)$. \hfill\textit{(1 pt)}
\item Calculer la covariance $\mathrm{Cov}(x,y)$ puis le coefficient de corrélation
linéaire $r$ (à $10^{-3}$ près). Un ajustement affine est-il justifié~? \hfill\textit{(1 pt)}
\item Déterminer une équation de la droite de régression de $y$ en $x$ par la méthode des
moindres carrés, puis estimer le montant épargné par un client ayant $8$ ans
d'ancienneté. \hfill\textit{(1 pt)}
\end{enumerate}

\smallskip
\noindent\textbf{Exercice 2.} \hfill\textit{(3 points)}\\
Un épargnant ouvre un compte à la microfinance avec un dépôt initial de \num{200000}
francs CFA. Chaque année, son capital est augmenté de $8\,\%$ d'intérêts, puis il verse
\num{100000} francs CFA supplémentaires. On note $u_n$ le capital (en francs CFA)
disponible au bout de $n$ années ; ainsi $u_0=\num{200000}$.
\begin{enumerate}[label=\arabic*)]
\item Justifier que, pour tout $n\in\N$, $u_{n+1}=1{,}08\,u_n+\num{100000}$. \hfill\textit{(0,5 pt)}
\item On pose $v_n=u_n+\num{1250000}$. Montrer que $(v_n)$ est géométrique ; préciser sa
raison et son premier terme. \hfill\textit{(1,25 pt)}
\item Exprimer $u_n$ en fonction de $n$. \hfill\textit{(0,75 pt)}
\item Déterminer le nombre d'années nécessaires pour que le capital dépasse
\num{1000000} de francs CFA. \hfill\textit{(0,5 pt)}
\end{enumerate}

\smallskip
\noindent\textbf{Exercice 3.} \hfill\textit{(4 points)}\\
Le solde (en millions de francs CFA) d'un fonds de roulement est modélisé par une fonction
$f$, dérivable sur $[0,+\infty[$, solution de l'équation différentielle
$(E):\ y'=-2y+6$ vérifiant la condition initiale $f(0)=1$ (le temps $t$ est en années).
\begin{enumerate}[label=\arabic*)]
\item Résoudre l'équation différentielle $(E)$ : déterminer l'ensemble de ses solutions
sur $\R$. \hfill\textit{(1,25 pt)}
\item Déterminer la solution particulière $f$ vérifiant $f(0)=1$. \hfill\textit{(0,75 pt)}
\item Étudier les variations de $f$ sur $[0,+\infty[$ et déterminer $\displaystyle\lim_{t\to+\infty}f(t)$ ;
interpréter ce résultat. \hfill\textit{(1 pt)}
\item Déterminer, par le calcul, l'instant $t$ (en années, à $10^{-2}$ près) à partir
duquel le solde dépasse \num{2,9} millions de francs CFA. \hfill\textit{(1 pt)}
\end{enumerate}

\smallskip
\noindent\textbf{Exercice 4.} \hfill\textit{(5 points)}\\
Le plan complexe est muni d'un repère orthonormé direct $(O\,;\,\vec u,\vec v)$.
\begin{enumerate}[label=\arabic*)]
\item Résoudre dans $\C$ l'équation $z^2-2z+2=0$. On note $z_A$ la solution de partie
imaginaire positive et $z_B$ l'autre. \hfill\textit{(1 pt)}
\item Écrire $z_A$ sous forme trigonométrique et en déduire $z_A^{\,4}$. \hfill\textit{(1 pt)}
\item On considère les points $A$, $B$, $C$ d'affixes respectives $z_A=1+i$, $z_B=1-i$
et $z_C=3+i$. Démontrer que le triangle $ABC$ est rectangle isocèle et préciser le sommet
de l'angle droit. \hfill\textit{(1,5 pt)}
\item Soit $r$ la rotation de centre $A$ et d'angle $\dfrac{\pi}{2}$. Donner l'écriture
complexe de $r$, puis montrer que $C$ est l'image de $B$ par $r$. \hfill\textit{(1,5 pt)}
\end{enumerate}

\subsection*{Partie B — Évaluation des compétences \hfill (5 points)}

\noindent\textit{Madame Ngo Bell, gérante d'une microfinance à Douala, te sollicite en
tant qu'élève de Terminale~C pour conseiller une cliente. \textbf{(i)} La cliente hésite
entre deux plans d'épargne sur $5$ ans pour un même dépôt initial de \num{500000} francs
CFA : le \emph{Plan A} rapporte un intérêt simple en ajoutant \num{60000} francs CFA
chaque année au capital ; le \emph{Plan B} rapporte un intérêt composé de $9\,\%$ par an.
\textbf{(ii)} La cliente veut savoir, pour le Plan B, au bout de combien d'années son
capital aura \emph{doublé}. \textbf{(iii)} Sur les dossiers de prêt traités par l'agence,
$4\,\%$ se révèlent insolvables ; un inspecteur prélève $12$ dossiers au hasard et signale
l'agence dès qu'il trouve au moins un dossier insolvable.}

\begin{enumerate}[label=\textbf{Tâche \arabic*.},leftmargin=2.4em]
\item Déterminer le plan qui rapporte le plus à la cliente au bout de $5$ ans, et de
combien (en francs CFA). \hfill\textit{(1,5 pt)}
\item Déterminer, pour le Plan B, le nombre d'années au bout duquel le capital aura
\emph{doublé}. \hfill\textit{(1,5 pt)}
\item Déterminer, à $10^{-2}$ près, la probabilité que l'inspecteur signale l'agence. \hfill\textit{(1,5 pt)}
\end{enumerate}
\par\smallskip{\footnotesize\itshape Présentation générale : 0,5 point.}

% ════════════════════════════════════════════════════
\corrige{du sujet 15}

\noindent\textbf{Partie A — Exercice 1.}
1) $\overline{x}=\dfrac{1+2+3+4+5}{5}=3$ ; $\overline{y}=\dfrac{7+11+14+18+20}{5}=\dfrac{70}{5}=14$.
$V(x)=\dfrac15\sum(x_i-\overline{x})^2=\dfrac{4+1+0+1+4}{5}=\dfrac{10}{5}=2$.
2) $\mathrm{Cov}(x,y)=\dfrac15\sum(x_i-\overline{x})(y_i-\overline{y})$. Les écarts
$(x_i-3)$ valent $-2,-1,0,1,2$ et $(y_i-14)$ valent $-7,-3,0,4,6$, donc
$\sum=(-2)(-7)+(-1)(-3)+0+4+12=14+3+0+4+12=33$ et $\mathrm{Cov}(x,y)=\dfrac{33}{5}=6{,}6$.
$V(y)=\dfrac15(49+9+0+16+36)=\dfrac{110}{5}=22$. Donc
$r=\dfrac{\mathrm{Cov}(x,y)}{\sqrt{V(x)\,V(y)}}=\dfrac{6{,}6}{\sqrt{2\times22}}=\dfrac{6{,}6}{\sqrt{44}}\approx\mathbf{0{,}995}$.
Comme $|r|$ est très proche de $1$, l'ajustement affine est tout à fait justifié.
3) Pente $a=\dfrac{\mathrm{Cov}(x,y)}{V(x)}=\dfrac{6{,}6}{2}=3{,}3$ ; ordonnée à l'origine
$b=\overline{y}-a\overline{x}=14-3{,}3\times3=14-9{,}9=4{,}1$. La droite de régression est
$y=3{,}3\,x+4{,}1$. Pour $x=8$ : $y=3{,}3\times8+4{,}1=26{,}4+4{,}1=30{,}5$, soit environ
\textbf{\num{3050000} francs CFA} ($30{,}5$ centaines de milliers).

\smallskip
\noindent\textbf{Exercice 2.}
1) Au bout d'une année, le capital $u_n$ produit $8\,\%$ d'intérêts (il devient
$1{,}08\,u_n$), puis l'épargnant verse \num{100000}\,FCFA : d'où
$u_{n+1}=1{,}08\,u_n+\num{100000}$.
2) $v_n=u_n+\num{1250000}$, donc $u_n=v_n-\num{1250000}$ et
$v_{n+1}=u_{n+1}+\num{1250000}=1{,}08\,u_n+\num{100000}+\num{1250000}
=1{,}08(v_n-\num{1250000})+\num{1350000}=1{,}08\,v_n-\num{1350000}+\num{1350000}=1{,}08\,v_n$.
$(v_n)$ est géométrique de raison $q=1{,}08$ et de premier terme
$v_0=u_0+\num{1250000}=\num{1450000}$.
3) $v_n=\num{1450000}\times1{,}08^{\,n}$, d'où
$u_n=\num{1450000}\times1{,}08^{\,n}-\num{1250000}$.
4) $u_n>\num{1000000}\iff\num{1450000}\times1{,}08^{\,n}>\num{2250000}
\iff1{,}08^{\,n}>\dfrac{2250}{1450}=\dfrac{45}{29}\approx1{,}5517$, soit
$n>\dfrac{\ln(45/29)}{\ln1{,}08}\approx5{,}71$. Donc à partir de la \textbf{6\textsuperscript{e} année}
($u_6\approx\num{1050968}$\,FCFA $>\num{1000000}$).

\smallskip
\noindent\textbf{Exercice 3.}
1) $(E):\ y'=-2y+6$. Les solutions de l'équation homogène $y'=-2y$ sont
$y_h(t)=C\mathrm{e}^{-2t}$, $C\in\R$. Une solution particulière constante $y_p=k$ vérifie
$0=-2k+6$, soit $k=3$. Les solutions de $(E)$ sont donc
$y(t)=C\mathrm{e}^{-2t}+3$, $C\in\R$.
2) $f(0)=1\Rightarrow C+3=1\Rightarrow C=-2$. Donc $f(t)=3-2\mathrm{e}^{-2t}$.
3) $f'(t)=-2\times(-2)\mathrm{e}^{-2t}=4\mathrm{e}^{-2t}>0$ : $f$ est strictement croissante
sur $[0,+\infty[$. Comme $\mathrm{e}^{-2t}\to0$, $\displaystyle\lim_{t\to+\infty}f(t)=3$ :
le solde croît et tend vers une valeur d'équilibre de $3$ millions de francs CFA.
4) $f(t)>2{,}9\iff3-2\mathrm{e}^{-2t}>2{,}9\iff2\mathrm{e}^{-2t}<0{,}1\iff\mathrm{e}^{-2t}<0{,}05
\iff-2t<\ln0{,}05\iff t>-\dfrac{\ln0{,}05}{2}=\dfrac{\ln20}{2}\approx\mathbf{1{,}50}$ an.
Le solde dépasse $2{,}9$ millions de francs CFA à partir de $t\approx1{,}50$ an.

\smallskip
\noindent\textbf{Exercice 4.}
1) $z^2-2z+2=0$ : $\Delta=4-8=-4=(2i)^2$, $z=\dfrac{2\pm2i}{2}=1\pm i$. Donc
$z_A=1+i$ et $z_B=1-i$.
2) $|z_A|=\sqrt{1^2+1^2}=\sqrt2$ et $\arg z_A=\dfrac{\pi}{4}$, donc
$z_A=\sqrt2\,\mathrm{e}^{i\pi/4}$. D'où
$z_A^{\,4}=(\sqrt2)^4\mathrm{e}^{i\pi}=4\times(-1)=-4$.
3) $\vec{AB}$ a pour affixe $z_B-z_A=-2i$, donc $AB=2$. $\vec{AC}$ a pour affixe
$z_C-z_A=(3+i)-(1+i)=2$, donc $AC=2$. Enfin $z_C-z_B=(3+i)-(1-i)=2+2i$, donc
$BC=|2+2i|=2\sqrt2$. On a $AB=AC=2$ et $AB^2+AC^2=4+4=8=BC^2$ : par la réciproque du
théorème de Pythagore, $ABC$ est rectangle en $A$ et isocèle. Le sommet de l'angle droit
est $\mathbf{A}$.
4) La rotation de centre $A$ d'affixe $z_A$ et d'angle $\dfrac\pi2$ a pour écriture
complexe $z'-z_A=\mathrm{e}^{i\pi/2}(z-z_A)$, soit $z'=i(z-z_A)+z_A=i\,z+(1-i)z_A$. Avec
$z_A=1+i$ : $(1-i)(1+i)=2$, donc $z'=i\,z+2$. Image de $B$ :
$r(z_B)=i(1-i)+2=(i+1)+2=3+i=z_C$. Donc $C=r(B)$. \emph{(Cohérent avec la question 3 :
$AB=AC$ et $\widehat{BAC}=\frac\pi2$.)}

\smallskip
\noindent\textbf{Partie B — Situation.}
\emph{Tâche 1.} \textbf{Plan A} (intérêt simple, $+\num{60000}$/an pendant $5$ ans) :
$A_5=\num{500000}+5\times\num{60000}=\num{800000}$\,FCFA. \textbf{Plan B} (intérêt composé
$9\,\%$) : $B_5=\num{500000}\times1{,}09^5$. Or $1{,}09^5\approx1{,}53862$, donc
$B_5\approx\num{769311}$\,FCFA. Comparaison : $A_5-B_5\approx\num{800000}-\num{769311}
=\num{30689}$\,FCFA. C'est le \textbf{Plan A} qui rapporte le plus, d'environ
\textbf{\num{30689} francs CFA}.

\emph{Tâche 2.} Plan B : $C_n=\num{500000}\times1{,}09^n$. Le capital double quand
$1{,}09^n\ge2$, soit $n\ge\dfrac{\ln2}{\ln1{,}09}\approx\dfrac{0{,}6931}{0{,}0862}\approx8{,}04$.
Donc le capital aura doublé au bout de \textbf{$9$ ans} ($1{,}09^9\approx2{,}1719$, alors
que $1{,}09^8\approx1{,}9926<2$).

\emph{Tâche 3.} Soit $Z$ le nombre de dossiers insolvables parmi $12$ :
$Z\hookrightarrow\mathcal{B}(12\,;\,0{,}04)$. L'inspecteur signale l'agence si $Z\ge1$ :
$P(Z\ge1)=1-(0{,}96)^{12}\approx1-0{,}6127=\mathbf{0{,}39}$, soit environ \SI{39}{\percent}.
